\documentclass[11pt,a4paper]{article}
\usepackage[utf8]{inputenc}
\usepackage[T1]{fontenc}
\usepackage[ngerman]{babel}
\usepackage{geometry}
\usepackage{xcolor}
\usepackage{booktabs}
\usepackage{longtable}
\usepackage{enumitem}
\usepackage{listings}
\usepackage{hyperref}
\geometry{left=24mm,right=24mm,top=22mm,bottom=24mm}
\hypersetup{colorlinks=true,linkcolor=blue!45!black,urlcolor=blue!55!black}
\setlist{nosep,leftmargin=*}
\setlength{\parindent}{0pt}
\setlength{\parskip}{0.55em}
\definecolor{codebg}{RGB}{246,248,250}
\definecolor{accent}{RGB}{24,91,110}
\lstset{
  basicstyle=\ttfamily\small,
  backgroundcolor=\color{codebg},
  frame=single,
  framerule=0.2pt,
  breaklines=true,
  columns=fullflexible,
  showstringspaces=false,
  keywordstyle=\color{accent}\bfseries
}

\begin{document}

\begin{titlepage}
  \centering
  {\Large Angewandte Modellierung und Systemsimulation\par}
  \vspace{1.6cm}
  {\Huge\bfseries Problem Set 12\par}
  \vspace{0.4cm}
  {\LARGE The Great Spreeland Logistics Sync\par}
  \vspace{1.0cm}
  {\large Protokollarchitektur, Dispatcher-Swarm und dynamische UI\par}
  \vfill
  \begin{tabular}{ll}
    Bearbeiter: & Nils Schön\\
    Semester: & Sommersemester 2026\\
    Stand: & 20. Juli 2026\\
    Implementierung: & Google ADK 2.3.0, MCP 1.28.1, Python 3.13
  \end{tabular}
  \vfill
  {\small Die Implementierung ist vollständig im Verzeichnis
  \texttt{spreeland\_dispatcher/} enthalten.}
\end{titlepage}

\tableofcontents
\newpage

\section{Ziel und Systemgrenze}

Ziel ist ein resilienter Logistikprozess für 2\,t Bio-Gurken von
Burg/Spreewald nach Cottbus. Der Dispatcher darf eine Route erst freigeben,
wenn aktuelle Infrastruktur- und Wetterdaten vorliegen. Er soll externe
Lieferanten anhand deklarierter Fähigkeiten finden, ein gültiges Angebot
auswählen und einen Checkout nur innerhalb einer nachweisbaren
Eigentümerautorisierung zulassen. Der Benutzer erhält währenddessen einen
Ereignisstrom und abschließend eine native Statuskarte.

Die sechs Protokolle lösen unterschiedliche Kopplungsprobleme. Sie sind keine
Konkurrenten, sondern bilden einen Stack:

\begin{longtable}{p{0.15\textwidth}p{0.24\textwidth}p{0.52\textwidth}}
\toprule
\textbf{Protokoll} & \textbf{Grenze} & \textbf{Aufgabe im Szenario}\\
\midrule
MCP & Agent--Tool/Daten & Discovery und Aufruf read-only Brückenwerkzeuge\\
A2A & Agent--Agent & Discovery von Wetter- und Lieferantenagenten, Aufgabenstatus\\
UCP & Commerce & Produkt, Angebot, Warenkorb, Checkout und Bestellstatus\\
AP2 & Autorisierung & Kryptographischer Beleg für Checkout und Zahlung\\
A2UI & UI-Beschreibung & Deklarative, frameworkunabhängige Lieferkarte\\
AG-UI & Agent--Frontend & Geordneter bidirektionaler Ereignisstrom\\
\bottomrule
\end{longtable}

\subsection{Architekturübersicht}

\begin{figure}[h]
\centering
\fbox{\begin{minipage}{0.92\textwidth}
\centering
\textbf{A2UI-Renderer / Benutzeroberfläche}\\
$\updownarrow$ \quad AG-UI-Ereignisse, Eingaben, Freigabe \quad $\updownarrow$\\[0.3em]
\textbf{Spreeland Dispatcher (Google ADK)}\\
\begin{tabular}{ccc}
$\swarrow$ MCP & $\downarrow$ A2A & $\searrow$ UCP + AP2\\
Stadt-Brückendaten & Wetter- und Lieferagenten & Händler, Trusted Surface, PSP
\end{tabular}
\end{minipage}}
\caption{Schichten und Vertrauensgrenzen des Systems}
\end{figure}

MCP hält Datenzugriff aus der Agentenlogik heraus. A2A behandelt entfernte
Agenten als opake Dienste: Der Dispatcher benötigt weder deren Prompts noch
deren Speicher. UCP und AP2 werden getrennt betrachtet: UCP beschreibt, was
gekauft und wie der Checkout abgewickelt wird; AP2 beweist, wer genau diesen
Checkout und die Zahlung autorisiert hat. A2UI beschreibt die Ansicht als
Daten, während AG-UI deren Transport und laufende Statusereignisse übernimmt.

\section{Exercise 1: Protocol Architecture Design}

\subsection{1. Infrastructure Discovery mit MCP}

Der Dispatch Agent ist MCP-Host/Client. Ein Infrastruktur-MCP-Server läuft
nahe an der städtischen PostgreSQL-Datenbank und veröffentlicht mindestens
\texttt{get\_bridge\_status(route?)} und
\texttt{get\_route\_capacity(route, vehicle\_weight\_t)}. Nach
\texttt{initialize} sendet der Client \texttt{tools/list}. Erst das Ergebnis
dieser Discovery wird dem Modell als erlaubtes Toolset angeboten. Für die
Routenentscheidung folgt ein \texttt{tools/call}; zurückgegeben werden Status,
Traglast, Wartungsfenster, Beobachtungszeit und Datenherkunft. Genau diese
List/Call-Trennung ist in der MCP-Toolspezifikation vorgesehen
\cite{mcp-tools}.

\textbf{Sicherheitsentwurf:}

\begin{itemize}
  \item Der Datenbanknutzer erhält ausschließlich \texttt{SELECT}-Rechte auf
        eine freigegebene View; Eingaben werden parametrisiert.
  \item Im Produktionsbetrieb wird Streamable HTTP mit TLS, kurzlebiger
        Workload-Identität und serverseitiger Autorisierung verwendet. Für die
        lokale Abgabe ist stdio geeigneter, weil keine Ports oder Secrets
        erforderlich sind.
  \item API- und Datenbank-Schlüssel stehen niemals in Prompt, Agent Card oder
        Quelltext. Sie kommen aus Umgebungsvariablen bzw. Secret Manager.
  \item Toolantworten sind nicht automatisch vertrauenswürdig. Der
        Dispatcher validiert Schema, Aktualität, Statuswert und Traglast
        deterministisch.
  \item Ein stdio-Server schreibt Protokollframes auf stdout und Logs nur auf
        stderr, damit der Transport nicht beschädigt wird.
\end{itemize}

Die Abgabe implementiert diesen Ablauf tatsächlich: Ein lokaler MCP-Prozess
wird gestartet, beide Werkzeuge werden entdeckt und
\texttt{get\_bridge\_status} wird aufgerufen. Die in-memory Datensätze ersetzen
nur die nicht zugängliche städtische Datenbank.

\subsection{2. Expert Consultation mit A2A}

Der fremde Weather-Predictor stellt eine Agent Card am standardisierten
Well-Known-Endpunkt \texttt{/.well-known/agent-card.json} bereit. Sie enthält
Identität, unterstützte Interfaces, Protokollversion, Streamingfähigkeit und
Skills. Der Dispatcher:

\begin{enumerate}
  \item lädt und validiert die Karte, prüft HTTPS, Herkunft und gegebenenfalls
        ihre JWS-Signatur;
  \item sucht einen Skill mit Tags wie \texttt{route-weather} und prüft
        Ein-/Ausgabe-Medientypen;
  \item wählt das erste kompatible Interface in der Präferenzreihenfolge;
  \item sendet Region, Route und Lieferfenster als A2A-Nachricht;
  \item verarbeitet bei deklariertem Streaming Status- und Artefakt-Updates,
        bis ein terminaler Taskzustand erreicht ist.
\end{enumerate}

Die A2A-Spezifikation definiert Agent Cards gerade für diese
Fähigkeits-Discovery und trennt Nachrichten, Tasks, Artefakte und Streaming
\cite{a2a-spec,a2a-discovery}. Statische Zugangsdaten gehören nicht in eine
Agent Card. In der lokalen Abgabe folgen drei JSON-Karten dem A2A-1.0-Schema;
ein deterministischer Adapter ersetzt den externen HTTPS-Endpunkt. Die
Auswahllogik bleibt beim Austausch des Adapters unverändert.

\subsection{3. Secure Fulfillment mit UCP und AP2}

Die Bestellung benötigt \emph{beide} Protokolle:

\begin{enumerate}
  \item \textbf{UCP-Discovery:} Der Lieferant veröffentlicht sein UCP-Profil.
        Der Dispatcher prüft Version, Checkout-Fähigkeit, Bezahlhandler und
        Signaturschlüssel.
  \item \textbf{Checkout:} Der Agent erzeugt eine Session mit exakt
        2\,000\,kg Bio-Gurken, Lieferziel Cottbus, Preis, Fracht, Steuer,
        Lieferfenster und einer Idempotency-ID. Der Händler gibt den
        finalisierten Checkout signiert zurück. UCP standardisiert Erzeugung,
        Aktualisierung und Abschluss der Checkout-Session
        \cite{ucp-overview,ucp-guide}.
  \item \textbf{AP2 Checkout Mandate:} Ein Hash bindet die Autorisierung an den
        händlersignierten, finalen Checkout. Der Eigentümer bestätigt auf
        einer Trusted Surface Händler, Ware, Menge, Gesamtpreis, Währung und
        Ablaufzeit. Damit kann der Händler prüfen, dass der Agent genau diesen
        Checkout kaufen darf.
  \item \textbf{AP2 Payment Mandate:} Eine zweite Autorisierung bindet das
        gewählte Zahlungsmittel und den exakten Betrag an denselben
        Checkout-Hash. Credential Provider, Netzwerk und Payment Processor
        verifizieren Signatur, Key Binding, Gültigkeit, Nonce und Constraints.
  \item \textbf{Receipt und Audit:} Signierte Checkout- und Payment-Receipts,
        Mandate-IDs, Checkout-Hash, UCP-Order-ID und Zeitstempel werden
        revisionssicher gespeichert. Wiederholungen mit derselben
        Idempotency-ID dürfen keine zweite Bestellung auslösen.
\end{enumerate}

AP2 unterscheidet Human-Present und Human-Not-Present. Für eine unmittelbare
Bestellung ist eine geschlossene, vom Benutzer bestätigte Autorisierung
passend. Bei autonomer Beschaffung muss der Eigentümer vorher offene Mandates
mit harten Grenzen signieren, beispielsweise \glqq maximal 4\,000 EUR,
ausschließlich zertifizierte Bio-Gurken, 2\,000 kg, zugelassene Händler,
gültig bis 20.07.2026 18:00 Uhr\grqq. Das geschlossene Mandate darf nur
erstellt werden, wenn jede Constraint deterministisch erfüllt ist. Die
aktuelle AP2-Spezifikation beschreibt Checkout und Payment Mandates sowie
deren kryptographische Bindung an den Checkout \cite{ap2-spec,ap2-auth}.

Die lokale Demo prüft Budget, kanonischen Warenkorb-Hash und Manipulation mit
HMAC. Das Feld
\texttt{COURSE\_DEMO\_ONLY\_NOT\_AP2\_CONFORMANT} verhindert eine
Verwechslung mit einer echten AP2-Signatur. Produktionsfreigabe ist erst mit
Trusted Surface und echten Mandates zulässig.

\subsection{4. Dynamic Visualization mit A2UI und AG-UI}

Der Agent sendet keine React- oder Flutter-Quelldatei. Stattdessen enthält
\texttt{a2ui\_delivery\_card.json} eine A2UI-v1.0-\texttt{createSurface}-
Nachricht mit dem offiziellen Basic Catalog, einer flachen Komponentenliste
und einem getrennten Datenmodell. Der Client besitzt einen vertrauenswürdigen
Katalog und mappt \texttt{Card}, \texttt{Column}, \texttt{Text} und
\texttt{Divider} auf seine nativen Widgets. Dadurch ist das Payload
frameworkunabhängig und enthält keinen ausführbaren Agentencode. Die A2UI
v1.0 Candidate Specification beschreibt
\texttt{createSurface}, \texttt{updateComponents},
\texttt{updateDataModel} und \texttt{deleteSurface}
\cite{a2ui-spec}.

AG-UI transportiert den Lauf als geordneten Eventstrom. Die Abgabe emittiert:

\begin{lstlisting}
RUN_STARTED
TEXT_MESSAGE_START
STEP_STARTED -> TEXT_MESSAGE_CONTENT -> STEP_FINISHED
STATE_SNAPSHOT
CUSTOM(name="a2ui", value=<A2UI message>)
TEXT_MESSAGE_END
RUN_FINISHED
\end{lstlisting}

AG-UI ist transportagnostisch; dieselben Events können als JSONL, SSE oder
WebSocket-Frames übertragen werden. Lifecycle-, Text-, Tool- und
State-Events folgen definierten Mustern \cite{agui-architecture,agui-events}.
Die Statusmeldungen enthalten operative Begründungen (Datenquelle, erfüllte
Constraints, Auswahlkriterium), aber keine verborgene Chain-of-Thought.

\section{Exercise 2: Implementing the Dispatcher Swarm}

\subsection{Abgabestruktur}

\begin{longtable}{p{0.38\textwidth}p{0.53\textwidth}}
\toprule
\textbf{Datei} & \textbf{Verantwortung}\\
\midrule
\texttt{agent.py} & Google-ADK-\texttt{root\_agent}, synchron definierter
MCP-\texttt{McpToolset}, A2A-Verhandlung als Function Tool\\
\texttt{bridge\_mcp\_server.py} & lokaler stdio-MCP-Server mit zwei
Infrastrukturwerkzeugen\\
\texttt{mcp\_client.py} & \texttt{initialize}, echte Tool-Discovery und
Tool-Aufruf\\
\texttt{agent\_cards/*.json} & A2A-1.0-Selbstauskünfte von Wetter- und
Lieferantenagenten\\
\texttt{a2a\_negotiation.py} & Skill-Matching, Angebotsabfrage,
Constraint-Prüfung und Auswahl\\
\texttt{payments.py} & klar markierter lokaler Manipulations- und Budgettest\\
\texttt{agui\_stream.py} & AG-UI-kompatible JSONL-Ereignisse\\
\texttt{a2ui\_delivery\_card.json} & einreichbares UI-Schema\\
\texttt{run\_spreeland\_dispatcher.py} & vollständiger Offline-Lauf\\
\texttt{tests/test\_spreeland\_}\newline\texttt{dispatcher.py} & fünf Integrations- und
Strukturtests\\
\bottomrule
\end{longtable}

Der ADK-Agent folgt dem vorgegebenen Starter:

\begin{lstlisting}[language=Python]
infrastructure_tools = McpToolset(
    connection_params=StdioConnectionParams(
        server_params=StdioServerParameters(
            command=sys.executable,
            args=[str(_server_script)]
        )
    )
)

root_agent = Agent(
    name="Spreeland_Dispatcher",
    model="gemini-2.5-flash",
    tools=[infrastructure_tools, negotiate_supplier_agents],
    instruction="Check bridges, consult agents, require AP2 approval."
)
\end{lstlisting}

Die ADK-Dokumentation empfiehlt genau diese Kombination aus
\texttt{McpToolset}, \texttt{StdioConnectionParams} und
\texttt{StdioServerParameters}; der Toolset kann synchron in
\texttt{agent.py} definiert und dadurch auch deployt werden
\cite{adk-mcp}.

\subsection{Deterministischer End-to-End-Lauf}

Der Offline-Lauf ist absichtlich unabhängig von einem Gemini-Key. Er zeigt
Protokollmechanik reproduzierbar, statt einen erfolgreichen Lauf durch eine
Cloudverbindung nur zu behaupten:

\begin{enumerate}
  \item MCP entdeckt
        \texttt{get\_bridge\_status} und
        \texttt{get\_route\_capacity}.
  \item Der Wetteragent wird über den Skill
        \texttt{route-weather} gefunden. Er bewertet Lübbenau als
        \texttt{LOW}, Vetschau als \texttt{MEDIUM} und Leipe als
        \texttt{HIGH}.
  \item Die Brücke SPB-014 ist offen, für 40\,t zugelassen und hat bei einem
        6,8-t-Fahrzeug eine ausreichende Reserve. Die Route via Lübbenau wird
        gewählt; die gesperrte Leipe-Brücke scheidet deterministisch aus.
  \item Zwei Lieferanten werden über
        \texttt{gherkin-wholesale} gefunden. Beide erfüllen Menge und
        Bio-Zertifizierung. Bio-Spree kostet
        $2\,000\cdot1{,}85+90=3\,790$ EUR; Lausitz kostet
        $2\,000\cdot1{,}74+380=3\,860$ EUR. Bio-Spree gewinnt nach
        niedrigstem Landed Cost.
  \item 3\,790 EUR liegen innerhalb des Demo-Budgets von 4\,000 EUR.
        Warenkorb-Hash und Demo-Signatur werden verifiziert; der Status weist
        weiterhin auf die notwendige echte AP2-Freigabe hin.
  \item Das Ergebnis wird als State Snapshot und A2UI-Karte gestreamt.
\end{enumerate}

\subsection{Ausführung und Verifikation}

\begin{lstlisting}
uv sync
uv run python run_spreeland_dispatcher.py
uv run python run_spreeland_dispatcher.py --format ag-ui
uv run pytest tests/test_spreeland_dispatcher.py -q
\end{lstlisting}

Der Testumfang prüft:

\begin{itemize}
  \item A2A-Agent-Card-Parsing und Gewinner der Verhandlung,
  \item echten MCP-Subprozess, Discovery und Tool-Aufruf,
  \item A2UI-Version, Root-Komponente, eindeutige IDs und Referenzen,
  \item Erkennung eines manipulierten Warenkorbs,
  \item vollständigen Dispatch-Flow und gewählte Brücke.
\end{itemize}

Für die optionale interaktive ADK-Weboberfläche wird
\texttt{GOOGLE\_API\_KEY} gesetzt und \texttt{uv run adk web .}
gestartet. Dieser Schritt benötigt als einziger einen externen Cloud-Key; der
prüfbare Kern der Einreichung funktioniert ohne ihn.

\section{Fehlerbehandlung, Emergenz und Skalierung}

Interoperabilität beseitigt nicht die komplexe Dynamik eines
Multi-Agenten-Systems. Sie macht sie jedoch beobachtbar und begrenzbar:

\begin{itemize}
  \item \textbf{Timeout und Circuit Breaker:} Wetter- und Händleragenten
        erhalten Fristen. Bei Ausfall gibt es keine erfundene Antwort, sondern
        einen degradierenden Status und einen alternativen Agenten.
  \item \textbf{Deterministische Gates:} Traglast, Zertifikat, Menge, Budget,
        Signaturen und Ablaufzeiten werden außerhalb des LLM geprüft.
  \item \textbf{Idempotenz:} Task-, Checkout- und Zahlungs-IDs verhindern, dass
        Retries doppelte Bestellungen erzeugen.
  \item \textbf{Feedback-Loops:} Jede neue Brücken- oder Wetterlage löst
        höchstens eine debouncte Neuplanung aus. Hysterese verhindert
        Route-Flapping bei kleinen Risikoänderungen.
  \item \textbf{Beobachtbarkeit:} AG-UI-Ereignisse korrelieren
        \texttt{threadId}, \texttt{runId}, Schritt, Tool und A2A-Task. Audits
        speichern Entscheidungen und Belege, nicht private Modellgedanken.
  \item \textbf{Least Privilege:} Der Dispatcher darf lesen und planen. Kaufen
        darf er erst nach einem gültigen AP2-Nachweis innerhalb der
        Constraints.
\end{itemize}

Damit können einzelne Agentenangebote variieren, ohne dass daraus unkontrollierte
Systemaktionen entstehen. Die Protokolle definieren die Leitplanken; harte
Policies und Signaturprüfungen bleiben deterministischer Anwendungscode.

\section{Ergebnis und Grenzen}

Die Aufgabenstellung ist vollständig in drei geforderten Artefaktklassen
umgesetzt: dieses Design-PDF, ausführbarer Python-Quelltext und ein
A2UI-JSON-Schema. Die Demo beweist MCP-Discovery, A2A-Skill-Matching,
Mehranbieter-Verhandlung, Routenwahl, Streaming sowie UI-Bindung.

Nicht vorgetäuscht werden externe Systeme, die das Aufgabenblatt nicht
bereitstellt: eine reale City-PostgreSQL-Instanz, öffentlich erreichbare A2A-
und UCP-Endpunkte, Eigentümer-Zahlungsschlüssel und ein AP2 Trusted Surface.
Die Austauschstellen sind in README und Code explizit markiert. Daher ist der
lokale Lauf vollständig reproduzierbar, während eine reale Zahlung
ordnungsgemäß blockiert bleibt, bis echte Autorität und Infrastruktur
konfiguriert sind.

\newpage
\begin{thebibliography}{9}

\bibitem{mcp-tools}
Model Context Protocol: \emph{Tools}.\\
\url{https://modelcontextprotocol.io/specification/draft/server/tools}

\bibitem{a2a-spec}
A2A Protocol Working Group: \emph{Agent2Agent Protocol Specification 1.0}.\\
\url{https://a2a-protocol.org/latest/specification/}

\bibitem{a2a-discovery}
A2A Protocol Working Group: \emph{Agent Discovery}.\\
\url{https://a2a-protocol.org/latest/topics/agent-discovery/}

\bibitem{ucp-overview}
Universal Commerce Protocol: \emph{Official Specification -- Overview}.\\
\url{https://ucp.dev/latest/specification/overview/}

\bibitem{ucp-guide}
Google for Developers: \emph{Universal Commerce Protocol Guide}.\\
\url{https://developers.google.com/merchant/ucp/guides}

\bibitem{ap2-spec}
Agent Payments Protocol: \emph{AP2 Specification}.\\
\url{https://ap2-protocol.org/ap2/specification/}

\bibitem{ap2-auth}
Agent Payments Protocol: \emph{Agent Authorization Framework}.\\
\url{https://ap2-protocol.org/ap2/agent_authorization/}

\bibitem{a2ui-spec}
A2UI Project: \emph{A2UI Protocol v1.0 Candidate}, aktualisiert am
8. Juni 2026.\\
\url{https://github.com/a2ui-project/a2ui/blob/main/specification/v1_0/docs/a2ui_protocol.md}

\bibitem{agui-architecture}
AG-UI: \emph{Core Architecture}.\\
\url{https://docs.ag-ui.com/concepts/architecture}

\bibitem{agui-events}
AG-UI: \emph{Events}.\\
\url{https://docs.ag-ui.com/concepts/events}

\bibitem{adk-mcp}
Google Agent Development Kit: \emph{MCP Tools}.\\
\url{https://adk.dev/tools-custom/mcp-tools/}

\end{thebibliography}

\end{document}
